As neural networks are deployed in critical settings, mechanistic interpretability aims to reverse-engineer their computations into analyzable units \cite{Elhage2021}. Sparse Autoencoders (SAEs) have proven effective at decomposing polysemantic MLP activations of language models into monosemantic features \cite{Cunningham2024, Bricken2023}. Applying this to Vision Transformers (ViTs) \cite{Dosovitskiy2021} raises two challenges: (i) interpreting sparse visual directions without built-in language alignment, and (ii) understanding how spatial evidence flows into the classification token. This project focuses on the SAE track of the problem: we train SAEs on ViT-B/16 MLP activations and use CLIP \cite{Radford2021} as an external evaluator for cross-modal concept labelling, providing a layer-wise empirical characterisation of learned visual features and their causal relevance. We further run the identical pipeline on a self-supervised DINOv2 backbone \cite{Caron2021} as a controlled comparison, quantifying how far a supervised-ViT SAE recipe transfers across pre-training paradigms. Direct attention-level tracing of patch-to-\texttt{[CLS]} aggregation is left for future circuit-level analysis.
